\documentclass{beamer}
\usetheme{Madrid}
\usecolortheme{default}

\title{Chess Provers in Lean}
\author{Lean Hackathon Team}
\date{May 22, 2026}

\begin{document}

\begin{frame}
  \titlepage
\end{frame}

\begin{frame}{Agenda}
\begin{itemize}
  \item Current System
  \vspace{0.8cm}
  \item Alternative System
\end{itemize}
\end{frame}

\section{What We Did}

\begin{frame}{Approach Followed}
\begin{itemize}
  \item Built a proof-aware chess validation system in Lean.
  \item Input FEN, parse LLM-proposed moves, verify legality in Lean.
  \item Success only when the validated sequence ends in checkmate.
  \item Design split: heuristic generation (LLM), formal acceptance (Lean rules + proofs).
\end{itemize}
\end{frame}

\begin{frame}{System Map: \texttt{Main.lean}}
\begin{itemize}
  \item Entry flow in \texttt{main}:
  \begin{enumerate}
    \item Parse input FEN using \texttt{fenToGameState}
    \item Call Python client: \texttt{script.py <fen>}
    \item Parse JSON response via \texttt{parsePythonOutputToMoves}
    \item Validate sequence with \texttt{state.isValidMoves}
    \item Execute \texttt{playMoves} and check \texttt{isAnyCheckmate}
  \end{enumerate}
  \item Output includes success/failure and board-state sequence for valid mates.
\end{itemize}
\end{frame}

\begin{frame}{Trusting Initial State: \texttt{Chess/FEN.lean}}
\begin{itemize}
  \item Example FEN: \texttt{r2k2nr/pp1b1Q1p/2n4b/3N4/3q4/3P4/PPP3PP/4RR1K w - - 1 0}
  \item \texttt{fenToGameState} has the following helpers: \\
  \texttt{charToPiece, expandFenDigit, parseRankToVector, parseBoardToVector, parseCastling, parseEnPassant.}
  \item Invalid FEN sets \texttt{valid := false} with a safe fallback board/state.
\end{itemize}
\end{frame}

\begin{frame}{LLM Output to Typed Moves: \texttt{Chess/MoveParser.lean}}
\begin{itemize}
  \item JSON schema represented by \texttt{JsonMove} and \texttt{JsonCheckmateSolution}.
  \item Conversion chain:
  \begin{itemize}
    \item string decoders: \texttt{stringToPiece}, \texttt{stringToColour}, \texttt{stringToPos}
    \item structural conversion: \texttt{convertMoves : List JsonMove -> List Move}
  \end{itemize}
  \item \texttt{parsePythonOutputToMoves} is the typed boundary from untrusted JSON to Lean moves.
\end{itemize}
\end{frame}

\begin{frame}{Model Core (1): Data Types and Pseudo-Legality}
\begin{itemize}
  \item In \texttt{Chess/Model.lean}: core types \texttt{GameState}, \texttt{Move}, \texttt{Board}, \texttt{Pos}.
  \item Piece-specific move semantics encoded in:
  \begin{itemize}
    \item \texttt{isValidPawnMove}, \texttt{isValidKnightMove}, \texttt{isValidBishopMove}, etc.
  \end{itemize}
  \item \texttt{isPseudoLegalMove} enforces geometry/path/ownership constraints,
  but not king-safety after move.
\end{itemize}
\end{frame}

\begin{frame}{Model Core (2): Check, Legality, and Play Pipeline}
\begin{itemize}
  \item \texttt{isInCheck}: determines whether king is attacked.
  \item \texttt{isLegalMove}: strengthens pseudo-legality with king safety.
  \item State-transition and sequence validators:
  \begin{itemize}
    \item \texttt{makeValidMove}
    \item \texttt{isValidMoves}
    \item \texttt{playMoves}
  \end{itemize}
  \item End condition functions:
  \begin{itemize}
    \item \texttt{isCheckmate}
    \item \texttt{isAnyCheckmate}
  \end{itemize}
\end{itemize}
\end{frame}

\begin{frame}{Proof Coverage: \texttt{Chess/LegalProofs.lean}}
\begin{itemize}
  \item Focus theorem: \texttt{legal\_implies\_kingSafeAfter}
  \item Meaning: if \texttt{isLegalMove s m = true}, then
  \texttt{isInCheck (s.makeMove m) s.turn = false}.
  \item This links boolean legality checks to a concrete king-safety guarantee.
\end{itemize}
\end{frame}

\begin{frame}{Proof Coverage: \texttt{Chess/Proofs.lean}}
\begin{itemize}
  \item \texttt{no\_friendly\_fire}:
  valid target squares never contain a same-color piece.
  \item \texttt{bishop\_stays\_on\_colour}:
  bishop moves preserve board-square parity (color class).
  \item Role: demonstrate local rule invariants and movement properties as machine-checked facts.
\end{itemize}
\end{frame}

\begin{frame}{LLM Client: How It Works}
\centering
\Huge Demo
\end{frame}

\section{What Could Have Been}

\begin{frame}{Alternative Approach from \texttt{Mal-Pat/Chess-Lean}}
\begin{itemize}
  \item Repo centers around \texttt{ChessLean/Basic.lean}.
  \item Architecture leans toward internal move generation:
  \begin{itemize}
    \item \texttt{generateAllPseudoLegalMoves}
    \item \texttt{IsPseudoLegalMove} as membership
    \item \texttt{IsLegalMove := pseudo-legal \&\& not-in-check}
  \end{itemize}
  \item Execution style: \texttt{playMove}/\texttt{play} over generated legal-space transitions.
\end{itemize}
\end{frame}

\section{Alternative Chess Problems We Could Solve}

\begin{frame}{Closing}
\begin{itemize}
  \item We implemented a practical hybrid: external generation, internal formal acceptance.
  \item Proofs already anchor key safety and movement invariants.
  \item Next scope can expand from checkmate solving to broader verified chess reasoning.
\end{itemize}
\end{frame}

\begin{frame}
\centering
\Huge Thank You
\end{frame}

\end{document}
